\section{Introduction}
\label{sec:introduction}

% Original scaffold TODO: motivation, gap, benchmark, and contributions.
Video understanding involves knowledge about the practices, objects, and social settings that a video depicts.
A model may recognize a procession or a meal yet miss the particular tradition that gives it meaning.
Cultural benchmarks therefore ask models to interpret content across communities and languages, rather than assume that performance on a general benchmark transfers uniformly to every cultural setting.
This concern motivates cultural image and video benchmarks
\citep{nayak-etal-2024-benchmarking,shafique-etal-2025-culturally,Singh_2026_CVPR}.

Cultural cues also appear together in ways that do not fit a single label for the whole video.
Consider a lion dance performed during a community event outside China.
The form, the participants, the event, and the language heard in the recording may have different cultural associations.
The setting can help explain the event without changing the identity of the performance.
Similarly, a recognizable textile may become part of a garment associated with another tradition.
These examples motivate cultural mixing as an evaluation setting that includes coexistence and recontextualization as well as adaptation and fusion.
% Illustrative examples, not claims about adjudicated benchmark items.

Existing work provides several parts of this evaluation problem.
Cultural question answering assesses knowledge about culturally situated content.
Mixed-context image evaluation examines whether a target remains recognizable when other cultural cues are present \citep{Kim_2026_CVPR}.
Video adds procedures, changing social roles, speech, music, and evidence that may recur at separated moments.
An assessment of cultural mixing should therefore establish both which cues are present and how the evidence supports their interpretation.
It should also distinguish recognizing two cues from understanding the relation between them.

We develop \benchmark{} around this distinction.
The basic unit is a culturally associated cue instance, such as a practice, material, sound, visible text, or event setting.
Each cue has a description, a component type, one or more supported cultural associations, and an evidence record.
The proposed core tasks ask models to recover these cues and identify the channels that support them.
A candidate relation task connects cues and their participating entities or contexts through a relation frame with explicit roles.
For example, a practice and the group performing it occupy different roles; two parallel cultural displays may instead have a symmetric relation.
Role assignment expresses the organization visible in a sample, without asserting the historical direction of cultural influence.

Our initial retrieval exploration focused on performing arts and media adaptations.
The intended scope also includes food, clothing and materials, ceremonies, craft, and other everyday practices.
\todo{Confirm the final domain coverage after audiovisual screening.}

\begin{itemize}
    \item We define cultural mixing in video to include coexistence, recontextualization, and adaptation.
    \item We represent each cue through its cultural associations, component type, and supporting evidence.
    \item We propose tasks for cue recognition and modality attribution, plus a candidate task for classifying relations and assigning roles. Modality and text-only controls test which inputs support each prediction.
\end{itemize}
